% GENERATED FILE. Edit canonical lessons, then run npm run generate:books.
% canonical-source-hash: fnv1a64:6521dfb0ea4c74bf
% canonical-lessons: FR-C475-jalousie, FR-C475-envier, FR-C475-rancune, FR-C475-rancunier, FR-C475-reconcilier

\chapter{Jalousie et rancune}
\label{ch:fr-jalousie-et-rancune}
\hlchaptermodality{475}

\begin{chapteropening}
\textbf{By the end of this chapter:} \emph{I can talk about jealousy, grudges and making up.}

Complete the last lesson of chapter 475: I can talk about jealousy, grudges and making up.
\end{chapteropening}

\section[la jalousie]{la jalousie — jealousy}

\label{lesson:FR-C475-jalousie}



\begin{quote}
Before the new one: say the French for funny, then the French for to tickle.
\end{quote}

\subsection*{You'll want to know: la jalousie}
\textbf{la jalousie} — \textquotedblleft{}jealousy\textquotedblright{}.

\begin{quote}
La jalousie a gâché leur amitié.
\end{quote}

\begin{quote}
\emph{Jealousy spoilt their friendship.}
\end{quote}

Built on \emph{jaloux}. Also a slatted blind on a window.

\begin{cousinweb}
\emph{la jalousie} is built on \emph{jaloux} (EtymDB), which goes back to Late Latin \emph{zēlōsus} \textquotedblleft{}full of zeal\textquotedblright{}, from \emph{zēlus}, a borrowing from Greek \emph{zêlos} \textquotedblleft{}zeal, rivalry\textquotedblright{}.
\end{cousinweb}

\begin{grammarlens}[title={What you've built}]
One more word for telling, explaining and describing at B1.
\end{grammarlens}

\subsection*{Your turn}
\begin{itemize}
\raggedright
  \item \emph{Say it:} \emph{la jalousie}
  \item \emph{Say it:} \emph{la jalousie}, once more
  \item \emph{Say it:} a sentence of your own with \emph{la jalousie}
  \item \emph{Recall:} say the French for funny, then the French for to tickle, then say \emph{la jalousie} again
\end{itemize}

\begin{tcolorbox}[breakable,colback=teal!4,colframe=teal!35!black,title={Before you move on}]
What does la jalousie mean? (\textquotedblleft{}Jealousy\textquotedblright{}.) Say it once more.
\end{tcolorbox}
\section[envier]{envier — to envy}

\label{lesson:FR-C475-envier}



\begin{quote}
Before the new one: say the French for to tickle, then the French for jealousy.
\end{quote}

\subsection*{You'll want to know: envier}
\textbf{envier} — \textquotedblleft{}to envy\textquotedblright{}.

\begin{quote}
J'envie mon frère : il voyage beaucoup.
\end{quote}

\begin{quote}
\emph{I envy my brother: he travels a lot.}
\end{quote}

\emph{Envier quelque chose à quelqu'un} — to envy someone something.

\begin{cousinweb}
\emph{envier} is built on \emph{envie} (EtymDB), inherited from Latin \emph{invidia} \textquotedblleft{}envy\textquotedblright{}, from \emph{invidēre} \textquotedblleft{}to look askance at\textquotedblright{}, built on \emph{vidēre} \textquotedblleft{}to see\textquotedblright{} (Lewis \& Short).
\end{cousinweb}

\begin{grammarlens}[title={What you've built}]
One more word for telling, explaining and describing at B1.
\end{grammarlens}

\subsection*{Your turn}
\begin{itemize}
\raggedright
  \item \emph{Say it:} \emph{envier}
  \item \emph{Say it:} \emph{envier}, once more
  \item \emph{Say it:} a sentence of your own with \emph{envier}
  \item \emph{Recall:} say the French for to tickle, then the French for jealousy, then say \emph{envier} again
\end{itemize}

\begin{tcolorbox}[breakable,colback=teal!4,colframe=teal!35!black,title={Before you move on}]
What does envier mean? (\textquotedblleft{}To envy\textquotedblright{}.) Say it once more.
\end{tcolorbox}
\section[la rancune]{la rancune — grudge, resentment}

\label{lesson:FR-C475-rancune}



\begin{quote}
Before the new one: say the French for jealousy, then the French for to envy.
\end{quote}

\subsection*{You'll want to know: la rancune}
\textbf{la rancune} — \textquotedblleft{}grudge, resentment\textquotedblright{}.

\begin{quote}
Je ne garde pas de rancune contre lui.
\end{quote}

\begin{quote}
\emph{I do not bear him a grudge.}
\end{quote}

\emph{Sans rancune !} — no hard feelings!

\begin{cousinweb}
\emph{la rancune} comes from Old French \emph{rancure}, which goes back to Latin \emph{rancor} \textquotedblleft{}a rancid smell, a grudge\textquotedblright{} (Lewis \& Short); its ending was later reshaped.
\end{cousinweb}

\begin{grammarlens}[title={What you've built}]
One more word for telling, explaining and describing at B1.
\end{grammarlens}

\subsection*{Your turn}
\begin{itemize}
\raggedright
  \item \emph{Say it:} \emph{la rancune}
  \item \emph{Say it:} \emph{la rancune}, once more
  \item \emph{Say it:} a sentence of your own with \emph{la rancune}
  \item \emph{Recall:} say the French for jealousy, then the French for to envy, then say \emph{la rancune} again
\end{itemize}

\begin{tcolorbox}[breakable,colback=teal!4,colframe=teal!35!black,title={Before you move on}]
What does la rancune mean? (\textquotedblleft{}Grudge, resentment\textquotedblright{}.) Say it once more.
\end{tcolorbox}
\section[rancunier]{rancunier — resentful, unforgiving}

\label{lesson:FR-C475-rancunier}



\begin{quote}
Before the new one: say the French for to envy, then the French for grudge.
\end{quote}

\subsection*{You'll want to know: rancunier}
\textbf{rancunier} — \textquotedblleft{}resentful, unforgiving\textquotedblright{}.

\begin{quote}
Ma voisine est rancunière : elle n'oublie rien.
\end{quote}

\begin{quote}
\emph{My neighbour is resentful: she forgets nothing.}
\end{quote}

Feminine \emph{rancunière}. Someone who keeps grudges.

\begin{cousinweb}
\emph{rancunier} is built on \emph{rancune} \textquotedblleft{}grudge\textquotedblright{}, which goes back to Latin \emph{rancor} \textquotedblleft{}a rancid smell, a grudge\textquotedblright{} (Lewis \& Short).
\end{cousinweb}

\begin{grammarlens}[title={What you've built}]
One more word for telling, explaining and describing at B1.
\end{grammarlens}

\subsection*{Your turn}
\begin{itemize}
\raggedright
  \item \emph{Say it:} \emph{rancunier}
  \item \emph{Say it:} \emph{rancunier}, once more
  \item \emph{Say it:} a sentence of your own with \emph{rancunier}
  \item \emph{Recall:} say the French for to envy, then the French for grudge, then say \emph{rancunier} again
\end{itemize}

\begin{tcolorbox}[breakable,colback=teal!4,colframe=teal!35!black,title={Before you move on}]
What does rancunier mean? (\textquotedblleft{}Resentful, unforgiving\textquotedblright{}.) Say it once more.
\end{tcolorbox}
\section[réconcilier]{réconcilier — to reconcile; (se réconcilier) to make up}

\label{lesson:FR-C475-reconcilier}



\begin{quote}
Before the new one: say the French for grudge, then the French for resentful.
\end{quote}

\subsection*{You'll want to know: réconcilier}
\textbf{réconcilier} — \textquotedblleft{}to reconcile; (se réconcilier) to make up\textquotedblright{}.

\begin{quote}
Après des mois de silence, ils se sont réconciliés.
\end{quote}

\begin{quote}
\emph{After months of silence, they made up.}
\end{quote}

Mostly reflexive: \emph{se réconcilier avec quelqu'un}.

\begin{cousinweb}
\emph{réconcilier} is a learned borrowing from Latin \emph{reconciliāre} \textquotedblleft{}to bring together again\textquotedblright{}, built on \emph{concilium} \textquotedblleft{}a gathering\textquotedblright{} (Lewis \& Short).
\end{cousinweb}

\begin{grammarlens}[title={What you've built}]
One more word for telling, explaining and describing at B1.
\end{grammarlens}

\subsection*{Your turn}
\begin{itemize}
\raggedright
  \item \emph{Say it:} \emph{réconcilier}
  \item \emph{Say it:} \emph{réconcilier}, once more
  \item \emph{Say it:} a sentence of your own with \emph{réconcilier}
  \item \emph{Recall:} say the French for grudge, then the French for resentful, then say \emph{réconcilier} again
\end{itemize}

\begin{tcolorbox}[breakable,colback=teal!4,colframe=teal!35!black,title={Before you move on}]
What does réconcilier mean? (\textquotedblleft{}To reconcile; (se réconcilier) to make up\textquotedblright{}.) Say it once more.
\end{tcolorbox}
